\paragraph{A proved radial subfamily.}
For either fixed rooted quartet shape, take the original finite-$M$ Moran
MSRC recursion with equal branch copy counts, an interior point-mass initial
count, neutral zero-switching dynamics, and zero lower-branch durations.
Condition on one of $0011$, $0101$, or $0110$, and vary only the root duration
$T$. The paired pattern has a compatible pair whose coalescence determines the
corresponding quartet topology $Q_j$. Neutral complement and taxon-exchange
symmetries give equal probabilities for the two nonselected resolutions.
Thus
\[
q_j(T)=d(T),\qquad q_k(T)=q_\ell(T)=(1-d(T))/2.
\]
At $T=0$, $q=(1/3,1/3,1/3)$; as $T\to\infty$, $d(T)\to1$ while the
conditioning event remains positive at every finite $T$. Continuity and the
intermediate-value theorem therefore imply that the open radial segment from
the simplex center to each corresponding vertex is attained.

This theorem is strictly one-dimensional. Exact pectinate calculations with
asymmetric lower branches move $q_1/(q_1+q_3)$ on both sides of $1/2$, so no
simple universal ordering obstruction is available. However, no theorem shows
that this ratio spans its full admissible interval at fixed $q_2$. Hence full
conditional simplex coverage remains unresolved.
